\documentclass[11pt]{article}
\usepackage[margin=1.1in]{geometry}
\usepackage{amsmath,amssymb,amsthm,booktabs,tabularx,array}
\usepackage{microtype,xcolor}
\usepackage[colorlinks=true,linkcolor=blue!50!black,citecolor=blue!50!black,urlcolor=blue!50!black]{hyperref}
\newtheorem{theorem}{Theorem}[section]
\newtheorem{lemma}[theorem]{Lemma}
\newtheorem{corollary}[theorem]{Corollary}
\newtheorem{proposition}[theorem]{Proposition}
\newcommand{\N}{\mathbb N}
\newcommand{\code}[1]{{\small\ttfamily #1}}
\setlength{\emergencystretch}{1.5em}
\title{Finite Parity Defects Force Collatz Returns:\\
An Explicit Near-Periodicity Certificate in Lean}
\author{M.~Sharpe\thanks{Developed with Codex (OpenAI), building on the author's existing Collatz formalization. No claim of priority in the mathematical literature is made. The Collatz conjecture remains open.}}
\date{September 5, 2026}
\begin{document}
\maketitle
\begin{abstract}
Write $T$ for the shortcut Collatz map and suppose $N+1\le2^K$.
We prove that at most $D$ disagreements between the parity itinerary of
$N$ and its $q$-step shift, over at least
$(2^{D+1}-1)(q+K)$ comparisons, force an exact $q$-step orbit return.
The return occurs at one of the explicit times
$(2^k-1)(q+K)$, $0\le k\le D$.
A parity prefix within $e$ letter edits of a periodic template has at
most $2e$ shift disagreements, giving a robust finite extension of the
exact prefix-power criterion. The proof combines a simple orbit-height
bound with parity congruence and disjoint geometric windows; it imposes
no restrictions on edit positions or odd-step density.
For positive $q$ the certificate bounds the entire orbit. For $q=2$
and a positive seed it proves that the orbit reaches one. The general
return, template, boundedness, and period-two results are checked in
Lean~4 using only standard foundational axioms. We also prove a sharper
universal height bound, reducing the geometric window-growth factor from
two to $8/5$, with no additional orbit hypothesis. The new horizons are
never longer than the original ones. Universal availability of these
finite certificates is not established.
\end{abstract}

\section{Definitions and the finite return theorem}
Use the shortcut map
\[
T(n)=\begin{cases}n/2,&n\equiv0\pmod2,\\(3n+1)/2,&n\equiv1\pmod2,\end{cases}
\qquad n\in\N=\{0,1,2,\ldots\},
\]
and write $b_t=T^t(N)\bmod2$. For a lag $q\ge1$ define the shift-defect set
\begin{equation}\label{eq:defects}
\mathcal D_{N,q}(L)=\{t\in\N: t<L,\ b_t\ne b_{t+q}\},\qquad
d_{N,q}(L)=|\mathcal D_{N,q}(L)|.
\end{equation}
Thus $L$ comparisons require the parity prefix of length $L+q$.
Choose $K\in\N$ with $N+1\le2^K$; the smallest such $K$ is
$\lceil\log_2(N+1)\rceil$. Define
\begin{equation}\label{eq:checkpoints}
A=q+K,\qquad s_0=0,\qquad s_{k+1}=2s_k+A.
\end{equation}
Induction gives $s_k=(2^k-1)A$.

\newpage
\begin{theorem}[Finite defect certificate]\label{thm:main}
Suppose $N+1\le2^K$, $q\ge1$, and
\begin{equation}\label{eq:threshold}
L\ge(2^{D+1}-1)(q+K),\qquad d_{N,q}(L)\le D.
\end{equation}
Then some $k\le D$ satisfies
\begin{equation}\label{eq:return}
T^{s_k+q}(N)=T^{s_k}(N).
\end{equation}
Consequently the tail from $s_k$ is periodic with period dividing $q$,
and the entire orbit is bounded.
\end{theorem}
The formal return theorem is
\code{ParityDefects.few\_\allowbreak{}defects\_\allowbreak{}force\_\allowbreak{}return}.
Its return conclusion is also valid for $q=0$; we assume $q\ge1$ in
the paper to obtain a nonzero period and boundedness. A return does not
by itself identify the cycle as the trivial Collatz cycle.

\section{A return or a nearby disagreement}
\begin{lemma}[Elementary height bound]\label{lem:height}
For every $N,u\in\N$,
\begin{equation}\label{eq:height}
T^u(N)+1\le2^u(N+1).
\end{equation}
In particular, if $N+1\le2^K$, then $T^u(N)<2^{u+K}$.
\end{lemma}
\begin{proof}
Both branches satisfy $T(n)+1\le2(n+1)$. Iterate this inequality and
use integrality for the strict upper bound.
\end{proof}
This deliberately coarse bound retains the effect of the added one.
It makes no assumption about which parity branch is taken.

\begin{lemma}[Local alternative]\label{lem:local}
Suppose $N+1\le2^K$. At every time $s$, either
$T^{s+q}(N)=T^s(N)$, or there is a shift disagreement at some
\begin{equation}\label{eq:local}
s\le t<2s+q+K.
\end{equation}
\end{lemma}
\begin{proof}
Suppose no disagreement occurs in that interval. Put $r=s+q+K$,
$x=T^s(N)$, and $y=T^{s+q}(N)$. The first $r$ parity letters of
$x$ and $y$ agree. The previously formalized parity congruence theorem
therefore gives $x\equiv y\pmod{2^r}$.
If $x\ne y$, their difference is a nonzero multiple of $2^r$, so
$2^r\le\max(x,y)$. But Lemma~\ref{lem:height} gives
\[
x<2^{s+K}\le2^r,\qquad y<2^{s+q+K}=2^r,
\]
a contradiction. Thus $x=y$.
\end{proof}
In Lean this uses \code{prefix\_\allowbreak{}power} from the existing
prefix-power development \cite{rung1}. That theorem packages the same
parity congruence argument for the $q$-periodic block of length $q+r$.
The new local alternative is \code{return\_\allowbreak{}or\_\allowbreak{}parity\_\allowbreak{}defect}.

\section{Disjoint windows count the defects}
\begin{proof}[Proof of Theorem~\ref{thm:main}]
Suppose all $D+1$ checkpoints $s_0,\ldots,s_D$ fail to return.
Lemma~\ref{lem:local} supplies a disagreement in each interval
\[
[s_k,2s_k+q+K)=[s_k,s_{k+1}),\qquad 0\le k\le D.
\]
These intervals are disjoint and contained in $[0,s_{D+1})\subseteq[0,L)$.
Hence there are at least $D+1$ distinct disagreements, contrary to
$d_{N,q}(L)\le D$.

An exact return at time $s$ gives
$T^{s+q+j}(N)=T^{s+j}(N)$ for every $j\ge0$, by determinism of the map.
Every later orbit value is therefore among the first $s+q$ values.
Their finite maximum bounds the whole orbit, including the initial
segment before the return.
\end{proof}
The formal counting lemma is
\code{ParityDefects.count\_\allowbreak{}ge\_\allowbreak{}of\_\allowbreak{}no\_\allowbreak{}checkpoint\_\allowbreak{}return}.
It proves the cardinality bound by inserting one new defect at each
step. Boundedness is supplied by \code{return\_\allowbreak{}bounds\_\allowbreak{}orbit},
which uses strong induction to reduce any time beyond $s+q$ by $q$.

\begin{corollary}[Necessary defect growth]\label{cor:unbounded}
If the orbit of $N$ is unbounded, then for every $q\ge1$, $D\ge0$,
and $K$ with $N+1\le2^K$,
\begin{equation}\label{eq:many}
d_{N,q}\bigl((2^{D+1}-1)(q+K)\bigr)>D.
\end{equation}
\end{corollary}
\begin{proof}
Otherwise Theorem~\ref{thm:main} bounds the orbit.
\end{proof}
For fixed $N,q,K$ this is a logarithmic necessary growth condition in
the observation horizon. It is not a positive-density assertion about
the disagreements. An orbit can fail this finite certificate while
having a very sparse defect set.

\section{Letter edits from a periodic template}
Let $w_t\in\{0,1\}$ be a template satisfying
$w_{t+q}=w_t$ for $0\le t<L$. Let
\begin{equation}\label{eq:errors}
E=\{t<L+q:b_t\ne w_t\}.
\end{equation}
Only this finite template segment is relevant.

\begin{lemma}\label{lem:edits}
$d_{N,q}(L)\le2|E|$.
\end{lemma}
\begin{proof}
If $b_t\ne b_{t+q}$ while $w_t=w_{t+q}$, at least one of $t$ and
$t+q$ is in $E$. Thus
\[
\mathcal D_{N,q}(L)\subseteq E\cup\{i-q:i\in E,\ i\ge q\}.
\]
Both sets on the right have cardinality at most $|E|$.
The Lean proof uses the image of all of $E$ under natural subtraction
by $q$; possible extra zero indices do not affect the upper bound.
\end{proof}

\begin{theorem}[Near-power certificate]\label{thm:template}
Suppose $N+1\le2^K$, $q\ge1$, the template is $q$-periodic on the
specified segment, and $|E|\le e$. If
\begin{equation}\label{eq:editbound}
L\ge(2^{2e+1}-1)(q+K),
\end{equation}
then $T^{s_k+q}(N)=T^{s_k}(N)$ at some $k\le2e$. In particular the
orbit is bounded.
\end{theorem}
\begin{proof}
Use Lemma~\ref{lem:edits} and Theorem~\ref{thm:main} with $D=2e$.
\end{proof}
The edit locations are unrestricted. There is no assumption of a
Sturmian slope, positive density gap, or bounded ideal correction.
The tradeoff is explicit: the sufficient horizon grows exponentially
with the number of edited letters.

\subsection*{Why the direct swap-cost argument was insufficient}
The earlier affine-word identity \cite{silver} says that swapping
adjacent $01$ and $10$, after prefix $u$ and before suffix $v$, changes
the correction by
\[
2^{|u|}3^{\operatorname{ones}(v)}.
\]
Taking $u$ to be $i$ zero letters and $v$ empty gives a single swap
with correction change $2^i$. A count of swaps alone cannot furnish a
uniform small correction bound. The new proof instead uses the exact
intervals between disagreements and does not estimate an accumulated
weighted correction. It therefore establishes a different, explicit
sufficient condition for approximate repetition.

\section{Period two and a worked certificate}
\begin{lemma}\label{lem:two}
If $n>0$ and $T^2(n)=n$, then $n=1$ or $n=2$.
\end{lemma}
\begin{proof}
Split by the parities of $n$ and $T(n)$. The even--even branch gives
$n/4=n$, hence $n=0$, excluded. The even--odd branch gives
$(3n+2)/4=n$, hence $n=2$. The odd--even branch gives
$(3n+1)/4=n$, hence $n=1$. The odd--odd branch gives
$(9n+5)/4=n$, which has no positive solution.
\end{proof}

\begin{corollary}\label{cor:one}
For a positive seed, either certificate with $q=2$ implies that the
orbit reaches one.
\end{corollary}
\begin{proof}
Positivity is preserved by $T$. Apply Lemma~\ref{lem:two} at the
certified return. If that value is two, the next value is one.
\end{proof}

For example, $N=3$, $K=2$, $q=2$ gives $A=4$. Its orbit begins
\[
3,5,8,4,2,1,2,1,\ldots,\qquad
b_0b_1b_2\cdots=11000101\cdots.
\]
The shift disagreements are at $t=0,1,3$. With $D=3$ the sufficient
comparison horizon is $L=(2^4-1)4=60$, and the checkpoints are
$0,4,12,28$. The first positive checkpoint already returns:
$T^4(3)=T^6(3)=2$.
The template $w_t=t\bmod2$ differs only at times $0$ and $3$.
For this template $e=2$, so \eqref{eq:editbound} gives $L=124$ and
uses $126$ parity letters. This illustrates the distinction between
shift disagreements and edited letters. These counts and thresholds
were also checked using Lean's kernel-checked \code{decide}.

\section{Sharper windows from the actual growth bound}\label{sec:sharp}
The original height estimate uses one bit per step. The shortcut map
permits a smaller uniform allowance while retaining the added-one correction.

\begin{lemma}\label{lem:sharpheight}
For every natural $N,u$,
\begin{equation}\label{eq:threehalves}
2^u\bigl(T^u(N)+1\bigr)\le3^u(N+1).
\end{equation}
Define $g(u)=\lceil3u/5\rceil=\lfloor(3u+4)/5\rfloor$.
If $N+1\le2^K$, then
\begin{equation}\label{eq:sharpheight}
T^u(N)<2^{g(u)+K}.
\end{equation}
Moreover, $g$ is nondecreasing and $g(u)\le u$.
\end{lemma}
\begin{proof}
Both branches satisfy $2(T(n)+1)\le3(n+1)$; induction gives
\eqref{eq:threehalves}. Write $u=5a+r$, $0\le r<5$. We have
$3^5=243\le256=2^8$ and the five elementary bounds
\[
\begin{array}{c|rrrrr}
r&0&1&2&3&4\\\hline
3^r&1&3&9&27&81\\
2^{r+g(r)}&1&4&16&32&128
\end{array}
\]
Since $g(u)=3a+g(r)$, these give
$3^u\le2^{8a+r+g(r)}=2^{u+g(u)}$.
Combine with \eqref{eq:threehalves} and $N+1\le2^K$, cancel $2^u$,
and use integrality. The remaining properties follow from the ceiling
formula on natural arguments.
\end{proof}

\begin{proposition}[Height-sensitive windows]\label{prop:heightwindow}
Suppose $H:\N\to\N$ is nondecreasing and $T^u(N)<2^{H(u)}$ for every $u$.
Set $h_0=0$ and $h_{k+1}=h_k+H(h_k+q)$. If $d_{N,q}(L)\le D$ and
$h_{D+1}\le L$, then $T^{h_k+q}(N)=T^{h_k}(N)$ at some $k\le D$.
\end{proposition}
\begin{proof}
Absence of defects in $[s,s+H(s+q))$ gives $H(s+q)$ matching parity
letters for $T^s(N)$ and $T^{s+q}(N)$. Both values are strictly below
$2^{H(s+q)}$, the first by monotonicity of $H$. They must therefore be
equal by parity congruence. Apply this local alternative to the disjoint
windows $[h_k,h_{k+1})$ and count as before.
\end{proof}
The envelope here is an explicit hypothesis. The next specialization
supplies it from the initial bit bound, without an additional orbit assumption.

\begin{theorem}[Sharper universal certificate]\label{thm:sharp}
Suppose $N+1\le2^K$. Define
\begin{equation}\label{eq:sharpcheckpoints}
\sigma_0=0,\qquad \sigma_{k+1}=\sigma_k+
\left\lceil\frac{3(\sigma_k+q)}5\right\rceil+K.
\end{equation}
If $d_{N,q}(L)\le D$ and $L\ge\sigma_{D+1}$, an exact $q$-step return
occurs at some $\sigma_k$, $k\le D$. The conclusions about boundedness
for $q>0$ and convergence for $q=2$, $N>0$, still apply. The template
certificate also holds with $L\ge\sigma_{2e+1}$.
\end{theorem}
\begin{proof}
Use $H(u)=g(u)+K$ in Proposition~\ref{prop:heightwindow}, justified by
Lemma~\ref{lem:sharpheight}. The earlier return and template lemmas then apply.
\end{proof}

\begin{proposition}[Comparison and geometric rate]\label{prop:rate}
For all natural $q,K,k$,
\begin{equation}\label{eq:neverlonger}
\sigma_k\le(2^k-1)(q+K).
\end{equation}
With $C=3q+5K+4$, the following division-free bound also holds:
\begin{equation}\label{eq:rate}
5^k(3\sigma_k+C)\le8^kC.
\end{equation}
In particular, $\sigma_k\le\frac C3((8/5)^k-1)$.
\end{proposition}
\begin{proof}
Since $g(u)\le u$, induction compares the two checkpoint recurrences
and proves \eqref{eq:neverlonger}. For the sharper rate, the ceiling gives
$5\sigma_{k+1}\le8\sigma_k+C$, hence
$5(3\sigma_{k+1}+C)\le8(3\sigma_k+C)$.
Multiply by $5^k$ and induct from $\sigma_0=0$ to obtain \eqref{eq:rate}.
The final real inequality is its rearrangement.
\end{proof}

These worked examples compare exact horizons for $q=2$. The new horizons
and defect counts were checked with Lean's \code{decide}; the old horizons
use the original formula. These are not new computation records.
\begin{center}
\begin{tabular}{rrrrr}
\toprule
$N$&$K$&$D$&Original horizon&Sharper horizon\\
\midrule
3&2&3&60&36\\
7&3&5&315&128\\
9&4&6&762&250\\
\bottomrule
\end{tabular}
\end{center}
For $N=3$, the new checkpoints are $0,4,10,20,36$, with defects at
$0,1,3$. For $N=7$, the five defects are at $1,4,5,7,9$; for $N=9$,
the six are at $1,3,6,7,9,11$. The improvement shortens the sufficient
certificate without proving that every orbit has few enough defects.

\section{Formalization and remaining coverage}
The original results are in \code{lean/Collatz/ParityDefects.lean};
the height-sensitive extension is in \code{lean/Collatz/DefectHeight.lean}.
The main interfaces, in namespace \code{Collatz.ParityDefects}, are:
\begin{center}
\begin{tabularx}{\textwidth}{@{}lX@{}}
\toprule
Declaration&Conclusion\\
\midrule
\code{few\_\allowbreak{}defects\_\allowbreak{}force\_\allowbreak{}return}&An explicit checkpoint has a $q$-step return.\\
\code{near\_\allowbreak{}power\_\allowbreak{}force\_\allowbreak{}return}&The same from finitely many template edits.\\
\code{near\_\allowbreak{}power\_\allowbreak{}bound\_\allowbreak{}orbit}&A positive-period certificate bounds the orbit.\\
\code{unbounded\_\allowbreak{}forces\_\allowbreak{}defects}&The necessary defect bound \eqref{eq:many}.\\
\code{few\_\allowbreak{}defects\_\allowbreak{}two\_\allowbreak{}reaches\_\allowbreak{}one}&A positive seed with the period-two certificate reaches one.\\
\code{near\_\allowbreak{}two\_\allowbreak{}power\_\allowbreak{}reaches\_\allowbreak{}one}&The period-two template-edit version.\\
\code{sharp\_\allowbreak{}few\_\allowbreak{}defects\_\allowbreak{}force\_\allowbreak{}return}&A return at the shorter universal horizon.\\
\code{sharpCheckpoint\_\allowbreak{}growth}&The certified $8/5$ geometric upper bound.\\
\bottomrule
\end{tabularx}
\end{center}
The modules build with Lean \code{v4.27.0-rc1}. The expanded command
\code{python3 analysis/audit\_\allowbreak{}theorems.py --build}
audits 44 selected declarations. All inspected transitive dependencies
are within \code{propext}, \code{Classical.choice}, and \code{Quot.sound}.
The new general modules contain no \code{sorry},
\code{native\_\allowbreak{}decide}, or added mathematical axiom.
The finite worked controls used ordinary \code{decide} in a temporary
Lean file. This verification is not independent kernel replay or a
fresh rebuild of every external dependency and large certificate.

The theorem excludes a concrete family of approximate repetitions from
unbounded trajectories, extending exact repetition without restricting
the placement of defects. The horizons are conservative and exponential;
directly checking a long certificate is not proposed as a faster
algorithm than observing the orbit return itself.
For arbitrary $q$, boundedness still leaves nontrivial cycles open.
For $q=2$ the cycle is trivial, but no theorem supplies a certificate
for every positive seed.

Indeed, every orbit that reaches one eventually alternates, so its
period-two shift has only finitely many disagreements. Taking a long
enough horizon then meets the certificate. Merely assuming universal
availability would rephrase the convergence problem. Further progress
requires independent restrictions on defect counts, or sharper local
height estimates that make the certificate applicable to a separately
justified class. The general Collatz coverage argument remains missing.

\begin{thebibliography}{9}
\raggedright
\bibitem{rung1} M.~Sharpe, \emph{Rung One of the Word-Complexity Ladder:
a Machine-Verified Exclusion of an Automatic Collatz Itinerary by
Mahler's Method in Degree One}, August 2026.
Repository file \code{paper/rung1.tex}; formal prefix-power development
\code{lean/Collatz/PrefixPower.lean}.
\bibitem{silver} M.~Sharpe, \emph{An Unconditional Sturmian Exclusion for
the Collatz Map at Slope $1/\sqrt2$}, September 2026.
Repository file \code{paper/silver.tex}; affine-word development
\code{lean/Collatz/WordAffine.lean}.
\bibitem{source} M.~Sharpe, \emph{Collatz at the Critical Line}, source
repository, \url{https://github.com/msharpe248/collatz}.
Research ledger: \code{analysis/RESEARCH\_\allowbreak{}STATUS.md}.
Verification metadata: \code{analysis/theorem\_\allowbreak{}audit.json}.
\end{thebibliography}
\end{document}
